\documentclass[a4,11pt]{aleph-notas}

% -- Paquetes adicionales
\usepackage{aleph-comandos}
\usepackage{booktabs}
\usepackage{pdflscape}
\hypersetup{
    urlcolor=colordef!70,
}

% -- Datos
\institucion{Facultad de Ciencias Exactas, Naturales y Ambientales}
\carrera{Catálogo STEM}
\asignatura{Matemática III}
\tema{Reto no. 2: Análisis masico de una pieza industrial}
\autor{Andrés Merino}
\fecha{Periodo 2026-1}

\logouno[0.16\textwidth]{Logos/logoPUCE_04_ac}
\definecolor{colortext}{HTML}{1957d1}
\definecolor{colordef}{HTML}{1957d1}
\fuente{montserrat}


\begin{document}

\encabezado

%%%%%%%%%%%%%%%%%%%%%%%%%%%%%%%%%%%%%%%%
\section{Resultado de aprendizaje}
%%%%%%%%%%%%%%%%%%%%%%%%%%%%%%%%%%%%%%%%

\begin{itemize}
    \item \textbf{RdA 3 — Criterio 3.3:} Emplea integrales múltiples para determinar áreas, superficies y volúmenes.
\end{itemize}

Específicamente, el estudiante será capaz de:
\begin{itemize}
    \item Describir una región tridimensional mediante inecuaciones y elegir el sistema de coordenadas adecuado.
    \item Calcular masa, centro de masa y momentos de inercia de un sólido con densidad variable mediante integrales triples.
    \item Aplicar cambios de variable (coordenadas cilíndricas y esféricas) para simplificar el cálculo de integrales.
    \item Interpretar los resultados en el contexto del diseño y análisis de piezas industriales.
    \item Verificar resultados analíticos mediante integración numérica en Python.
\end{itemize}


%%%%%%%%%%%%%%%%%%%%%%%%%%%%%%%%%%%%%%%%
\section{Descripción}
%%%%%%%%%%%%%%%%%%%%%%%%%%%%%%%%%%%%%%%%

En ingeniería mecánica, el análisis másico de una pieza —masa total, centro de masa y momentos de inercia— es fundamental para el diseño de mecanismos, la predicción del comportamiento dinámico y la selección de materiales. Cuando la densidad del material no es uniforme (por ejemplo, por gradientes térmicos, aleaciones o porosidad), este análisis requiere integrales triples sobre regiones tridimensionales complejas.

\subsection*{Pregunta esencial}
\begin{itemize}[leftmargin=*]
    \item ¿Cómo permiten las integrales triples calcular propiedades másicas de objetos reales con geometría y densidad no uniformes, y qué ventaja aporta elegir el sistema de coordenadas correcto?
\end{itemize}

\subsection*{Reto}

Tu desafío es seleccionar una pieza industrial con simetría de revolución o simetría esférica (rodamiento, buje, cono de transmisión, cúpula, etc.), modelar su geometría como una región $E \subset \mathbb{R}^3$ definida por inecuaciones, proponer una función de densidad variable $\rho(x,y,z)$ justificada físicamente y calcular analítica y numéricamente sus propiedades másicas.

Debes entregar:
\begin{itemize}[leftmargin=*]
    \item Un minitutorial en video (máx.\ 8 minutos) con la explicación conceptual y el recorrido del análisis.
    \item Un informe en PDF redactado en \LaTeX\ con el desarrollo matemático y los resultados.
    \item Un notebook de Jupyter o Google Colab reproducible con todo el código y las visualizaciones.
\end{itemize}

\subsection*{Preguntas guía}
\begin{itemize}[leftmargin=*]
    \item ¿Cómo se describe la región $E$ mediante inecuaciones en coordenadas cartesianas, cilíndricas o esféricas?
    \item ¿Qué consideración física justifica que la densidad varíe dentro de la pieza?
    \item ¿Por qué el cambio a coordenadas cilíndricas o esféricas simplifica las integrales en piezas con simetría de revolución?
    \item ¿Qué relación hay entre la posición del centro de masa y la simetría geométrica de la pieza?
    \item ¿Cómo afecta la distribución de masa a los momentos de inercia respecto a los distintos ejes?
    \item ¿Qué tan cercanos son los resultados analíticos y numéricos? ¿Dónde puede originarse la diferencia?
\end{itemize}

\subsection*{Actividades guía}
\begin{enumerate}[leftmargin=*, label={{\arabic*.}}]
    \item Seleccionar una pieza industrial real con simetría geométrica reconocible (por ejemplo: cilindro macizo con cavidad cónica, semiesfera hueca, toro simplificado). Incluir una imagen o esquema acotado con dimensiones en el informe.
    \item Definir la región $E$ que ocupa la pieza mediante inecuaciones en al menos un sistema de coordenadas. Verificar que los límites de integración son consistentes con las dimensiones reales de la pieza.
    \item Proponer una función de densidad $\rho(x,y,z)$ no constante que refleje alguna característica física del material (por ejemplo, densidad mayor en la parte más gruesa, gradiente radial por enfriamiento, etc.). Justificar la elección.
    \item Calcular analíticamente la {masa total}:
    \[
        m = \iiint_E \rho(x,y,z)\, dV,
    \]
    mostrando el cambio de variable utilizado y todos los pasos del cálculo.
    \item Calcular analíticamente las {coordenadas del centro de masa} $(\bar{x}, \bar{y}, \bar{z})$:
    \[
        \bar{x} = \frac{1}{m}\iiint_E x\,\rho\, dV, \quad
        \bar{y} = \frac{1}{m}\iiint_E y\,\rho\, dV, \quad
        \bar{z} = \frac{1}{m}\iiint_E z\,\rho\, dV.
    \]
    Interpretar el resultado en función de la simetría de la pieza.
    \item Calcular analíticamente al menos dos {momentos de inercia} respecto a los ejes coordenados, por ejemplo:
    \[
        I_z = \iiint_E (x^2+y^2)\,\rho\, dV.
    \]
    \item Verificar numéricamente los resultados de masa y centro de masa usando \texttt{scipy.integrate} (integración de Monte Carlo o cuadratura adaptativa) y comparar con los valores analíticos, reportando el error relativo.
    \item Generar una visualización 3D de la región $E$ con \texttt{matplotlib} o \texttt{plotly}, indicando la posición del centro de masa sobre la figura.
    \item Analizar cómo cambia el centro de masa si se duplica la densidad en una subregión específica de la pieza. Presentar este análisis como un experimento numérico.
    \item Redactar conclusiones sobre la influencia de la geometría y la densidad variable en las propiedades másicas, y reflexionar sobre las simplificaciones del modelo.
\end{enumerate}

\subsection*{Recursos}
\begin{itemize}[leftmargin=*]
    \item \href{https://docs.scipy.org/doc/scipy/reference/integrate.html}{Módulo \texttt{scipy.integrate} — integración numérica en Python}.
    \item \href{https://matplotlib.org/stable/gallery/mplot3d/index.html}{Galería de gráficos 3D con Matplotlib}.
    \item \href{https://plotly.com/python/3d-mesh/}{Visualización de sólidos 3D con Plotly}.
    \item \href{https://numpy.org/doc/stable/reference/random/generated/numpy.random.uniform.html}{Integración de Monte Carlo con NumPy}.
    \item Complementar con al menos dos fuentes académicas (libro o artículo) citadas en formato APA.
\end{itemize}


%%%%%%%%%%%%%%%%%%%%%%%%%%%%%%%%%%%%%%%%
\section{Producto}
%%%%%%%%%%%%%%%%%%%%%%%%%%%%%%%%%%%%%%%%

\subsection{Minitutorial explicativo}
Un video de máximo 8 minutos donde se explique:
\begin{itemize}[leftmargin=*]
    \item Qué es una integral triple y cómo se usa para calcular masa en un sólido con densidad variable.
    \item Cómo se describe la región $E$ y por qué se eligió el sistema de coordenadas utilizado.
    \item Qué representa el centro de masa y cómo se interpreta su posición en la pieza real.
    \item Un ejemplo concreto: dado un punto de la pieza, cómo contribuye a la integral y qué implica para el diseño.
\end{itemize}

\subsection{Informe de aplicación}
Un documento en PDF que contenga:
\begin{enumerate}[leftmargin=*, label={\textbf{\arabic*.}}]
    \item \textbf{Introducción:} descripción de la pieza, imagen o esquema, dimensiones y justificación de la elección.
    \item \textbf{Modelo geométrico:} definición de la región $E$, límites de integración y sistema de coordenadas elegido.
    \item \textbf{Función de densidad:} expresión de $\rho$ y justificación física.
    \item \textbf{Cálculo de la masa:} desarrollo analítico completo con cambio de variable.
    \item \textbf{Centro de masa:} cálculo analítico e interpretación geométrica.
    \item \textbf{Momentos de inercia:} cálculo de al menos dos momentos e interpretación.
    \item \textbf{Verificación numérica:} código, resultados y comparación con el valor analítico (error relativo).
    \item \textbf{Experimento de sensibilidad:} análisis del efecto de modificar la densidad en una subregión.
    \item \textbf{Conclusiones y limitaciones:} hallazgos principales y restricciones del modelo.
    \item \textbf{Bibliografía} en formato APA.
\end{enumerate}

\subsection{Entrega técnica}
\begin{itemize}[leftmargin=*]
    \item El informe debe ser elaborado en la plantilla \LaTeX\ institucional disponible en Overleaf: \href{https://www.overleaf.com/read/vhdczzsymytd#3bf1e4}{Formato Tareas PUCE}.
    \item Incluir el enlace al video con permisos de acceso activos.
    \item Incluir el enlace al notebook (Google Colab o repositorio) con celdas ejecutadas y visibles.
    \item Si los enlaces no permiten visualizar el material, la actividad no podrá ser calificada y se asignará \textbf{0}.
\end{itemize}


%%%%%%%%%%%%%%%%%%%%%%%%%%%%%%%%%%%%%%%%
\section{Rúbrica de calificación}
%%%%%%%%%%%%%%%%%%%%%%%%%%%%%%%%%%%%%%%%

En caso de presentar un video con voz en off, rostro no visible o fuera del rango de tiempo establecido, se reducirá un nivel a cada punto de la rúbrica.

\begin{landscape}
\begin{center}\footnotesize
\begin{tabular}{p{5cm}p{4cm}p{4cm}p{4cm}p{4cm}}
\toprule
\textbf{Indicador} & \textbf{Excelente} & \textbf{Bueno} & \textbf{Aceptable} & \textbf{Insuficiente} \\ \midrule
\textbf{Modelo geométrico y densidad (10 pts)}
    & Define correctamente la región $E$ con límites consistentes, elige el sistema de coordenadas adecuado y propone una densidad variable con justificación física clara (10)
    & Región y densidad correctas con justificación incompleta o sistema de coordenadas subóptimo (8)
    & Región con errores menores en los límites o densidad constante sin justificación (6)
    & Región incorrecta o densidad no definida (0) \\ \midrule
\textbf{Cálculo de masa y centro de masa (15 pts)}
    & Calcula analíticamente la masa y el centro de masa con cambio de variable correcto, mostrando todos los pasos e interpretando el resultado (15)
    & Cálculo correcto con pasos intermedios incompletos o interpretación parcial (12)
    & Cálculo de la masa correcto pero centro de masa con errores o sin interpretación (9)
    & Errores conceptuales en las integrales o resultados ausentes (0) \\ \midrule
\textbf{Momentos de inercia (10 pts)}
    & Calcula correctamente al menos dos momentos de inercia e interpreta su significado físico en la pieza (10)
    & Un momento correcto con interpretación parcial o segundo momento con errores menores (8)
    & Un solo momento calculado o interpretación ausente (6)
    & Cálculo incorrecto o ausente (0) \\ \midrule
\textbf{Verificación numérica y experimento (10 pts)}
    & Verifica masa y centro de masa numéricamente con error relativo reportado, y presenta el experimento de sensibilidad con análisis de resultados (10)
    & Verificación correcta sin experimento de sensibilidad, o experimento sin análisis (8)
    & Verificación con errores en la implementación o experimento sin conclusiones (6)
    & Verificación ausente o código no reproducible (0) \\ \midrule
\textbf{Presentación en video (5 pts)}
    & Exposición clara y fluida; rostro visible durante toda la presentación (5)
    & Exposición clara con ligeras interrupciones; rostro visible (4)
    & Exposición poco fluida o lectura excesiva; rostro visible parcial (3)
    & Voz en off o rostro no visible / exposición ininteligible (0) \\ \midrule
\textbf{Gestión del tiempo (5 pts)}
    & 8:00 $\pm$ 30 s (5)
    & 8:00 $\pm$ 60 s (4)
    & 8:00 $\pm$ 90 s (3)
    & Fuera de $\pm$ 90 s (0) \\ \bottomrule
\end{tabular}
\end{center}
\end{landscape}

\end{document}
